\section{Authorization Inconsistencies in SAGA}

This chapter examines how the studied reference implementation interprets contact
policy at its enforcement points. We distinguish the documented selection of
a rule from the interpretation of the selected budget. The analysis concerns
the fixed-policy setting of Chapter~3 and is based on source inspection;
runtime propagation to the three security boundaries requires separate
execution evidence. The findings concern the inspected historical revision,
not the current upstream status.

% Evidence baseline: source inspected at commit 12a0373258a5cd3c49d332e4f3af55d2be019c35.
% No protocol implementation or verification model was changed for this chapter.

\subsection{Intended Contact-Policy Semantics}

A contact rulebook is represented as an ordered list of rules, each carrying
a \texttt{pattern} and an integer \texttt{budget}. A pattern matches an agent
identifier, or AID. For contact from $I$ to $R$, the relevant object is $R$'s
rulebook evaluated against $I$'s AID: permission has a direction. Its evaluation
does not distinguish individual application commands or tool actions.
Throughout this chapter, \emph{policy budget} denotes the
\texttt{budget} field of a contact rule, not the communication quota associated
with an access-control token.

In \texttt{contact\_policy.py}, \texttt{check\_rulebook()} checks rule fields
and admits integer budgets from $-1$ upwards. The \texttt{match()} function
first checks AID syntax and then tests patterns with \texttt{fnmatch.fnmatch()}.
Its documentation states that the returned budget belongs to the most
specific matching pattern. The function \texttt{aid\_specificity()} supplies
the comparison score, using weighted contributions from literals and wildcard
constructs in the user and agent-name components. We use this numerical
ordering as the documented selection criterion, without assuming that it is
a general ordering by containment between pattern languages.

For a well-formed fixed rulebook with a unique highest-scoring match, let
$b^*$ denote the budget attached to that rule. For no match, take $b^*=0$,
consistent with the matcher's default result. The policy baseline is
\[
  \mathsf{Allow}_{\mathrm{spec}}(P_R,I)
  \Longleftrightarrow b^*>0.
\]
The positivity convention comes from the Provider's \texttt{/access} handler,
which rejects both negative and zero results. Thus the baseline combines a
documented rule-selection contract with the Provider's explicit interpretation
of the result. It does not assert a separate published policy language or
resolve conflicts between equally specific rules. A malformed AID produces a
separate error result and is outside the valid-identity cases considered here.

A positive policy decision is only permission to proceed through this check.
The Provider additionally uses resource availability and remaining contact
allowance, and later protocol stages impose their own checks. Accordingly,
this baseline specifies contact eligibility rather than sufficient conditions
for token issuance or application-message acceptance.

% Evidence: saga/common/contact_policy.py, check_rulebook(), match() docstring,
% aid_specificity(), and pattern_specificity_component().
% Evidence: saga/provider/provider.py, access() negative/zero branches.

\subsection{Rule-Specificity Selection Flaw}

The implementation of \texttt{match()} does not maintain the state needed by
its documented selection rule. It initializes \texttt{best\_pattern} to
\texttt{None} and the return budget to zero. For each matching rule, it compares
that rule's specificity with \texttt{aid\_specificity(best\_pattern)} and
updates the budget when the comparison succeeds. It never updates
\texttt{best\_pattern}. Consequently, later rules are compared with the initial
sentinel, not with the most specific rule already encountered.

This is more precise than saying that specificity is ignored: the comparison
is executed, but its reference value remains wrong. The score of
\texttt{None} is $-1$, whereas the scores computed for well-formed patterns
are nonnegative, including zero for the universal wildcard. Therefore every
matching rule handled normally by the scorer passes the comparison. For a
valid AID and a well-formed rule list, the returned budget is the budget of
the last matching rule encountered; with no match it remains zero.

The order involved is the rulebook's list order, not the order of keys within
each rule dictionary. The registration path carries the list in the
\texttt{contact\_rulebook} field, the Provider stores that field, and the Agent
loads it into its local state. None of these paths sorts the rules by
specificity. The matcher's ordinary list traversal consequently preserves
the distinction between earlier and later rules.

Consider two rules that both match the same valid identity $I$. Let $r_s$
have a strictly higher specificity score and budget $-1$, and let $r_g$ have
a lower score and budget $10$. This is an abstract policy example; it does
not assume that $I$ can edit $R$'s configuration. The documented selection is
$r_s$ in either ordering. Inspection of the loop yields the following
different implementation results:
\begin{center}
\begin{tabular}{lll}
\hline
Rule order & Intended budget & Returned budget \\
\hline
$[r_s,r_g]$ & $-1$ (deny) & $10$ \\
$[r_g,r_s]$ & $-1$ (deny) & $-1$ \\
\hline
\end{tabular}
\end{center}

To compare Boolean decisions, define
\[
  \mathsf{Allow}_{\mathrm{impl}}(P_R,I)
  \Longleftrightarrow \mathtt{match}(P_R,I)>0.
\]
This is the actual matcher result interpreted by the Provider's policy
check, not the success of the entire endpoint. For the first ordering above,
\[
  \mathsf{Allow}_{\mathrm{impl}}(P_R,I)
  \neq \mathsf{Allow}_{\mathrm{spec}}(P_R,I).
\]
The same logical rule set thus has an order-dependent implementation result
despite an unchanged unique maximum under the documented semantics. Equal
scores are unnecessary for this disagreement. The observation establishes a
policy-evaluation mismatch; it does not, by itself, establish any of the
credential or message events defined in Chapter~3.

% Evidence: saga/common/contact_policy.py, match() loop and aid_specificity(None).
% Evidence: saga/user/user.py, agent_register() rulebook input and application.
% Evidence: saga/provider/provider.py, registration storage and access() rulebook lookup.
% Evidence: saga/agent.py, Agent initialization of self.contact_rulebook.
% The two-rule table is a static derivation, not a reported execution experiment.

\subsection{Budget-Semantics Inconsistency}

The two enforcement points also differ in how they interpret an identical
matcher result. In \texttt{provider.py}, \texttt{access()} returns a denial for
\texttt{budget < 0} and separately for \texttt{budget == 0}. Only a positive
value proceeds to the resource-allocation part of the handler. Writing
$\mathsf{ProviderAllow}$ for this policy check alone gives
\[
  \mathsf{ProviderAllow}(b) \Longleftrightarrow b>0.
\]
This predicate excludes the endpoint's identity, registration, and resource
checks; a true value is not yet a Provider grant.

In \texttt{agent.py}, the incoming connection handler evaluates the local
rulebook against the initiating AID using \texttt{match()}. It rejects only
when the returned value is negative. The budget is not tested for equality
to zero at this point. For this particular precheck,
\[
  \mathsf{AgentPrecheckAllow}(b) \Longleftrightarrow b\geq0.
\]
Here ``allow'' means only that this conditional does not reject the
connection. It is neither a statement that a token is generated nor a
statement that an application message is accepted.
\begin{center}
\begin{tabular}{lll}
\hline
Policy result & Provider policy check & Agent precheck \\
\hline
$b<0$ & Reject & Reject \\
$b=0$ & Reject & Continue \\
$b>0$ & Continue & Continue \\
\hline
\end{tabular}
\end{center}

The difference includes an explicit zero-budget match and the zero returned
when no rule matches. At either source of zero, the Provider treats the
result as insufficient permission, while the Agent's precheck does not stop
the connection. However, the ordinary fresh-contact path obtains material
through the Provider first. Its zero-budget denial remains effective at that
boundary. The receiver-side difference therefore does not independently
establish a complete unauthorized execution.

Several distinct quantities must be kept separate when assessing this
behavior. The configured rule budget is an input to policy evaluation. The
Provider's \texttt{counter} stores a remaining allowance for each initiating
AID within a target Agent's record. A first allocation initializes the entry
to the selected budget minus one; subsequent allocations decrement that
entry, with a floor of zero. An exhausted entry blocks further allocation.
Despite its name, this is remaining contact allowance, not an accumulated
count of accepted application messages. Allocation consumes allowance before
the later credential exchange has necessarily completed.

Token-use state is maintained separately. The \texttt{generate\_token()}
function sets the token's \texttt{communication\_quota} from the configured
\texttt{Q\_MAX} value, not from the matched contact-rule budget. Token validation checks the
receiver's active-token state, expiration, remaining token quota, and
recipient access-control key. The conversation loop decrements token quota
and also has a separate \texttt{MAX\_QUERIES} guard. These checks concern token
use and local conversation processing. Neither their presence nor their
success establishes that the correct policy rule was selected, and the
zero-budget comparison cannot be interpreted as a token-quota exhaustion
test.

% Evidence: saga/provider/provider.py, access() sign checks and counter update.
% Evidence: saga/agent.py, handle_i_agent_connection(), generate_token(),
% token_is_valid(), receive_conversation(), and MAX_QUERIES.
% Evidence: saga/config.py, Q_MAX.

\subsection{Cross-Layer Authorization Inconsistency}

The findings concern two different relationships. The rule-selection flaw is
a \emph{specification-to-implementation inconsistency}: the matcher can return
a budget that disagrees with the documented most-specific-rule contract.
The zero-budget difference is an \emph{enforcement-to-enforcement
inconsistency}: two consumers of the same policy result apply different
predicates. In the latter case,
\[
  \mathsf{ProviderAllow}(0)
  \neq \mathsf{AgentPrecheckAllow}(0).
\]
The second disagreement does not depend on overlapping rules. Conversely,
the first can change the sign of the returned budget without involving zero.
They should not be assigned the same consequence merely because both concern
authorization.

The Provider and receiving Agent both invoke the common matcher. The Agent's
incoming policy check therefore supplies an additional enforcement point,
but it does not independently recover the documented rule selection. With
the same ordered rulebook and identity, both invocations inherit the same
selection error. Agreement between two implementations of a check is
insufficient when both rely on the same incorrect decision procedure.
At zero, even their interpretation of that shared result differs.

These observations concern the preservation of policy meaning along the
pipeline defined in Chapter~3. The Provider's successful response supplies
contact material, including one public one-time key and its signature, the
target Agent's certificate and public access-control material, and the target
user's certificate. It removes the rulebook and counter from that response.
This response is not an access-control token; the receiving Agent generates
the token at a later boundary.

The incoming handler performs its policy precheck before credential handling.
Subsequent identity, Provider-stamp, signature, and one-time-key checks can
still stop a new credential exchange. Their purpose differs from choosing the
applicable contact rule. In particular, the Provider stamp authenticates an
Agent card constructed at registration, rather than a decision selecting a
rule for this contact. Later, \texttt{token\_is\_valid()} checks credential
state without re-evaluating rule specificity. Thus successful cryptographic
checks cannot be substituted for the missing semantic agreement.

Cryptographic correctness does not repair an incorrect authorization
decision. Even assuming that signatures, certificates, token checks, and the
authenticated transport work as intended, those mechanisms do not establish
that a credential was made available in accordance with the user's policy.
This is a separation of obligations, not a claim that every policy mismatch
survives all later checks or that the entire protocol has been exercised.

% Evidence: saga/provider/provider.py, registration card/stamp and access() response.
% Evidence: saga/agent.py, incoming policy precheck, credential-handling branches,
% generate_token(), and token_is_valid().

\subsection{Security Consequences}

For the rule-selection flaw, source inspection directly establishes that a
documented denial can become a positive implementation decision. Because
\texttt{access()} uses that value for its policy check, the mismatch changes
whether execution proceeds towards allocation. The successful allocation
branch returns contact material. Conditional on the endpoint's other checks
and ordinary resource requirements succeeding, the mismatch can therefore
reach the Provider-grant boundary. This is a conditional control-flow
conclusion, not a report of an experimentally observed grant.

In the terminology of Chapter~3, that boundary would constitute a Level~1
violation when
\[
  \neg\mathsf{Allow}_{\mathrm{spec}}(P_R,I)
  \land \mathsf{ProviderGrant}(I,R).
\]
An incorrect positive matcher result alone is not this event: the material
must actually be released. The distinction matters because a policy mismatch
can coexist with an unsuccessful endpoint request.

A Level~2 consequence would additionally require
\[
  \neg\mathsf{Allow}_{\mathrm{spec}}(P_R,I)
  \land \mathsf{TokenIssued}(R,I,T).
\]
The shared matcher means that the receiver's policy precheck need not correct
the earlier decision, making propagation relevant to investigate. Nevertheless,
token generation occurs after additional checks and depends on valid protocol
state. The static mismatch neither satisfies those checks by itself nor
records a token being released. It is therefore insufficient as standalone
evidence of credential issuance.

Likewise, the Level~3 condition
\[
  \neg\mathsf{Allow}_{\mathrm{spec}}(P_R,I)
  \land \mathsf{MessageAccepted}(R,I,T,m)
\]
requires the receiving application boundary to be reached with an accepted
message and valid credential state. Entering a credential-handling branch,
generating a token, and accepting a protected message are distinct
observations. This chapter establishes no execution exhibiting this final
condition.

For the zero-budget discrepancy, the supported conclusion is narrower. The
Agent precheck disagrees with the baseline denial at zero, but the Provider's
own decision is unchanged: it rejects that value. The discrepancy alone
therefore does not produce a Provider grant on the fresh-contact path.
Neither token issuance nor message acceptance follows from the receiver's
failure to reject at this single check. Without additional execution
evidence, this finding is a latent authorization inconsistency, not an
independently confirmed access violation.

The two findings thus establish concrete policy-evaluation and enforcement
disagreements, with different implications for the intermediate boundaries.
They show why the authorization-soundness obligations of Chapter~3 cannot be
assumed from the implementation's cryptographic checks. An implementation
mismatch is not automatically an end-to-end exploit. Chapter~5 examines
matcher-level evidence separately from the modeled consequences of an assumed
authorization divergence at material release, token issuance, and message
acceptance.

% Evidence: saga/provider/provider.py, access() policy rejection, allocation, and response.
% Evidence: saga/agent.py, incoming precheck, token generation, and message validation.
% Scope: source-level mismatch and conditional propagation; no runtime exploit result.
